\chapter{Optionals}

If wick has a single reason to exist, it is the optional. Everything else ---
the static types, the ban on truthiness, the compile-time engine calls --- is
supporting structure around one claim: that ``this value might be absent''
should be a fact the compiler knows and the programmer must discharge, not a
runtime surprise that surfaces three functions away from its cause. This
chapter explains optionals fully, and tells the story of the bug that made
them non-negotiable.

\section{The problem, in Lua}

Start with the failure, because the feature is shaped exactly to its
contours. Kora Night keeps a high score in a save file. The Lua build loaded
it like this:

\begin{lstlisting}[language=luatwin]
local best = tonumber(lt.load_save("koranight_best") or "0")
\end{lstlisting}

\noindent
Every clause is reasonable. `lt.load_save` returns the saved string, or `nil`
if there is no save; the `or "0"` supplies a default for the `nil` case; and
`tonumber` parses the string to a number. It reads correctly, it passed
review, and it shipped.

Then a player force-quit the game during a save. A save is a write to disk,
and a write interrupted at the wrong instant leaves the file created but
empty --- zero bytes. On the next launch, `lt.load_save` opened that file,
read it successfully, and returned not `nil` but the empty string `""`. Now
walk the expression again with `""` in hand:

\begin{itemize}
\item `lt.load_save(...)` is `""` --- a real string, not `nil`.
\item `"" or "0"` is `""`, because in Lua the empty string is
\emph{truthy}. The fallback never fires.
\item `tonumber("")` is `nil`, because the empty string is not a number.
\item `local best = nil`, and every later use of `best` --- comparisons,
arithmetic, drawing the number --- now operates on `nil`.
\end{itemize}

\noindent
The game crashed on the first frame. Not at the save, not at the load, but
downstream, where `nil` finally met an operation that could not tolerate it.
The stack trace pointed at a line that was completely innocent. Four separate
Lua decisions had to line up for this --- truthy empty strings, a load that
returns `""` rather than `nil` for an empty file, `nil`-punning parse failure,
and no static checking to catch the `nil` before it spread --- and all four
are defensible in isolation. Together they turned an interrupted disk write
into a first-frame crash that reading the code would never reveal.

\section{The type that makes it unwritable}

wick's answer is a type. An operation that can fail to produce a value does
not return that value; it returns an \emph{optional}, written `T?`, which is
either a `T` or `nil`. And the compiler will not let you use a `T?` as a `T`.

\begin{lstlisting}[language=wick]
let s = lt.load_save("hi")       // s: str?
lt.print(s, 4, 4, 1, 1, 1, 1)    // COMPILE ERROR: expected str, got str?
\end{lstlisting}

\noindent
The second line does not compile. The error is \texttt{expected str, got
str?}, and it appears the instant you write the line, on your machine, not on
a player's. There is no path by which a `str?` reaches `lt.print`, which wants
a `str`. To get from `str?` to `str` you must state, in the program, what
happens when the value is `nil`. wick gives you two ways to do that.

\section{Unwrapping with \texttt{??}}

The nil-coalesce operator `??` unwraps an optional by supplying a default for
the `nil` case. `a ?? b` evaluates to `a` when `a` is non-`nil`, and to `b`
otherwise; the result type is the non-optional base type.

\begin{lstlisting}[language=wick]
let text = s ?? "no save yet"    // text: str
let n = num("maybe") ?? 0        // n: num  (num parses str -> num?)
\end{lstlisting}

\noindent
The default must match the base type: `s ?? "..."` is fine because `s` is
`str?` and the default is `str`; a mismatched default reports
\texttt{'??' default must be str} (or whatever the base type is). The right
side is evaluated lazily, only when the left is `nil`, so it is safe to put
work there. And `??` requires a genuine optional on its left --- using it on a
value that can never be `nil` is the error \texttt{'??' left side must be an
optional}, because there is nothing to coalesce.

Now return to the war story. The wick translation of that fatal Lua line is:

\begin{lstlisting}[language=wick]
let best = num(lt.load_save("koranight_wick_best") ?? "") ?? 0
\end{lstlisting}

\noindent
Trace it with the zero-byte file that crashed the Lua build.
`lt.load_save(...)` returns `str?`. Whether it is `nil` (no file) or `""`
(empty file) makes no difference to what follows, because the first `??`
turns \emph{both} into a `str`: `nil` becomes `""`, and `""` is already `""`.
Then `num("")` is `num?` --- and it is `nil`, because the empty string is not a
number. The second `??` catches that and yields `0`. There is no value of the
save file, present or absent, well-formed or empty, that produces anything but
a clean `num`. The bug is not fixed; it is \emph{unwritable}. Delete either
`??` and the program does not compile.

That is the difference worth internalising. In Lua the safe version and the
buggy version look almost identical, and the compiler accepts both; the
difference surfaces only when a specific rare input arrives in production. In
wick the buggy version --- the one missing a `??` --- is a compile error, and
the safe version is the only one that builds. The type system has moved the
discovery of the bug from a player's first frame to the author's first
save.

\section{Narrowing with \texttt{if x != nil}}

The second way to handle an optional is to \emph{narrow} it: test it against
`nil`, and inside the block where the test succeeded, use it as a plain `T`.

\begin{lstlisting}[language=wick]
if s != nil {
  // inside this block, s has type str, not str?
  lt.print(s, 4, 4, 1, 1, 1, 1)
}
\end{lstlisting}

\noindent
The comparison `s != nil` does double duty: it is a `bool` condition, and it
tells the type checker that within the guarded block `s` cannot be `nil`, so
its type there is `str`. This is the same move a reader makes intuitively ---
``we already checked it is not nil, so here it is a string'' --- made formal
and enforced.

Narrowing is the right tool when you want to do several things with the value
and a default would be awkward, or when the `nil` case has its own behaviour:

\begin{lstlisting}[language=wick]
let saved = lt.load_save("name")
if saved != nil {
  lt.print("WELCOME BACK " + saved, 4, 4, 1, 1, 1, 1)
} else {
  lt.print("NEW PLAYER", 4, 4, 1, 1, 1, 1)
}
\end{lstlisting}

\paragraph{The limits of narrowing.} Narrowing in v0.1 works on \emph{local}
variables. A module global is not narrowed by an `if`, because other code ---
another function, a later frame --- could change it between the test and the
use; to handle a global's optional, use `??`, or copy it into a local first
and narrow the local. Assigning a new optional value to a narrowed variable
inside the block \emph{un-narrows} it: the guarantee held only for the value
that was tested. And narrowing does not thread through `and`/`or` chains in
v0.1, which is why the compiler asks you to parenthesise combined conditions
(Chapter~3); when in doubt, narrow with a nested `if` or copy to a local. Each
of these is a known v0.1 boundary, listed in Chapter~8, not a bug.

\section{What returns an optional}

Optionals are not something you sprinkle on; they arrive from the operations
that can genuinely fail, and the compiler propagates them from there. The
sources in v0.1 are:

\begin{center}
\begin{tabular}{@{}llp{0.40\linewidth}@{}}
\toprule
\textbf{Expression} & \textbf{Type} & \textbf{\texttt{nil} when} \\
\midrule
\texttt{lt.load\_save(k)} & \texttt{str?} & no readable save \\
\texttt{num(s)} & \texttt{num?} & \texttt{s} is not a number (incl.\ \texttt{""}) \\
\texttt{pop(xs)} & \texttt{T?} & the list is empty \\
\texttt{m[key]} & \texttt{T?} & the key is absent \\
\texttt{env(name)} & \texttt{str?} & the variable is unset \\
\bottomrule
\end{tabular}
\end{center}

\noindent
Each is a place where, in another language, you would get a sentinel you might
forget to check: an empty string, a false, a `nil`, a special ``not found''
value. wick makes the failure a type, so forgetting to check it is not
possible --- the forgotten check is a compile error. Map lookup is the purest
example: `m[key]` is \emph{always} `T?`, even when you are sure the key is
present, because the type system cannot know that you are sure. The missing-key
case is forced into view at the call site, where `??` disposes of it in three
characters:

\begin{lstlisting}[language=wick]
let hp = stats["health"] ?? 100
\end{lstlisting}

\section{Optionals in the type system}

A few rules tie optionals into the rest of the language; they were stated in
Chapter~3 and are worth collecting here from the optional's point of view.

\begin{itemize}
\item \textbf{Widening is free.} A `T` may be used wherever a `T?` is wanted;
going from ``definitely present'' to ``possibly absent'' loses no safety.
\item \textbf{Unwrapping is required.} A `T?` may not be used where a `T` is
wanted. `??` and narrowing are the only ways across, and both make you handle
`nil`.
\item \textbf{Optionals do not nest.} There is no `T??`. An optional is a `T`
or `nil`, full stop.
\item \textbf{\texttt{nil} needs a type.} A bare `nil` has no inferable type;
`let x: str? = nil` supplies one. `nil` is meaningful only as the empty case
of some specific optional.
\item \textbf{Comparing a non-optional to \texttt{nil} is an error.} If `x` is
a plain `str`, then `x != nil` is dead code --- `x` can never be `nil` --- and
the compiler says \texttt{comparing non-optional to nil} rather than let you
write a check that can never fail.
\end{itemize}

\section{The discipline in practice}

In day-to-day wick you stop writing the defensive scaffolding that a dynamic
language demands, because the compiler tracks what the scaffolding was
guarding. You do not wrap every load and parse in an `if x ~= nil` chain; you
write `??` once where the value enters, and from there on the value is a plain
`T` that cannot betray you. The optional is not overhead. It is the removal of
overhead --- the manual, forgettable, review-it-and-still-miss-it overhead of
remembering which values might be `nil` --- replaced by a compiler that
remembers for you and refuses to forget.

The zero-byte save is one bug. The class it belongs to --- a fallible operation
whose failure is representable as an ordinary value and therefore skippable ---
is enormous, and it is the class wick was built to close. Every `T?` in a wick
program is a small proof that one member of that class has been handled.
